%! Inverses of Exponential Functions — Unit 2, Lesson 2.10 — Warm-Up
\documentclass[10pt]{article}
\usepackage{precalculus-article}
\usepackage{precalculus-boxes}
\newcommand{\TallMath}[1]{$\displaystyle #1\rule[-1.4em]{0pt}{3.2em}$}

\begin{document}

\pageheader{Unit 2, Lesson 2.10}{Warm-Up}
\namedateperiod

\vspace{0.15in}

\begin{enumerate}[leftmargin=1.8em, itemsep=14pt]

\item Evaluate each expression. Then answer the question below.

\vspace{6pt}
\begin{center}
\renewcommand{\arraystretch}{2.0}
\begin{tabular}{@{} c c c c @{}}
  $\log_2 8 = $ \blank{2cm} & \quad & $2^3 = $ \blank{2cm} & \\
\end{tabular}
\end{center}

\vspace{4pt}
What do you notice about the two results?

\writeline

\vspace{4pt}

\item Complete each expression. Then describe the pattern you see.

\vspace{6pt}
\begin{center}
\renewcommand{\arraystretch}{2.0}
\begin{tabular}{@{} c c c c @{}}
  $\log_3\!\left(3^5\right) = $ \blank{2cm} & \quad & $3^{\log_3 7} = $ \blank{2cm} & \\
\end{tabular}
\end{center}

\vspace{4pt}
What pattern do these results suggest about the relationship between $\log_3$ and $3^x$?

\writeline

\vspace{4pt}

\item The function $g(x) = 2^x$ passes through the point $(3,\, 8)$. What point must the
      \emph{inverse} of $g$ pass through? Explain how you know.

\vspace{4pt}
Point on the inverse: \blank{3cm}

\vspace{4pt}
\writeline

\end{enumerate}

\end{document}
